% DOCUMENTO 1 - RESUMEN ACADÉMICO (social_media.tex)
\documentclass{article}
\usepackage[spanish]{babel}
\usepackage{graphicx}
\usepackage{hyperref}
\usepackage{geometry}
\geometry{a4paper, margin=2cm}

\begin{document}

\section*{\centering Curso: Estrategias de Preservación de Valor con Criptoactivos\\Universidad Real de la Cámara de Comercio}

\subsection*{Objetivos Generales}
\begin{itemize}
    \item Comprender los fundamentos tecnológicos de blockchain y su rol en protección patrimonial
    \item Dominar metodologías para adquisición segura de activos digitales
    \item Analizar proyectos mediante evaluación fundamental y métricas en cadena
    \item Desarrollar estrategias defensivas combinando análisis técnico y económico
\end{itemize}

\subsection*{Público Objetivo}
Inversores y profesionales financieros que busquen alternativas seguras en entornos económicos volátiles. No se requieren conocimientos previos.

\subsection*{Contenido Destacado}
El curso integra teoría y práctica comenzando con los principios de blockchain, su evolución histórica y mecanismos de consenso. Analizamos Bitcoin como reserva de valor y Ethereum para contratos autoejecutables.

La primera semana cubre preservación de valor: funcionamiento de monedas estables, análisis de capitalización de mercado y evaluación de documentos técnicos. La segunda semana aplica análisis técnico complementario identificando puntos estratégicos de operación.

\subsection*{Instructor}
Daniel Velarde Kubber - Especialista en:
\begin{itemize}
    \item Estrategias defensivas con criptoactivos desde 2019
    \item Análisis fundamental de proyectos blockchain
    \item Desarrollo de sistemas de monitoreo en cadena
\end{itemize}

\subsection*{Logística}
\begin{itemize}
    \item Duración: 12 sesiones (24 horas)
    \item Modalidad: Presencial/Virtual
    \item Fecha: Octubre 2024
\end{itemize}

\end{document}